%==============================================================================
% APÊNDICE A — GLOSSÁRIO DE TERMOS TÉCNICOS
%==============================================================================
\chapter{Glossário de Termos Técnicos}
\label{apendice:glossario}

% Definições de entradas do glossário
\newglossaryentry{supervisionado}{
  name=Aprendizado Supervisionado,
  description={Paradigma de aprendizado de máquina no qual o modelo é
    treinado utilizando pares de entrada-saída rotulados. É o método
    mais comum para classificação de imagens odontológicas.}
}

\newglossaryentry{naosupervisionado}{
  name=Aprendizado Não Supervisionado,
  description={Paradigma onde o modelo identifica padrões em dados não
    rotulados, como agrupamento de tipos de lesões.}
}

\newglossaryentry{reforco}{
  name=Aprendizado por Reforço,
  description={Paradigma onde um agente aprende a tomar decisões através
    de interações com o ambiente, recebendo recompensas por ações corretas.}
}

\newglossaryentry{deeplearning}{
  name=Deep Learning,
  description={Subcampo do aprendizado de máquina que utiliza redes
    neurais com múltiplas camadas para aprender representações hierárquicas
    de dados, como imagens ou textos.}
}

\newglossaryentry{transferlearning}{
  name=Transfer Learning,
  description={Técnica que reutiliza um modelo pré-treinado em uma grande
    base de dados como ponto de partida para uma tarefa correlata com
    menos dados disponíveis.}
}

\newglossaryentry{cnn}{
  name=CNN (Rede Neural Convolucional),
  description={Arquitetura de rede neural profunda especializada no
    processamento de dados com grade topológica, como imagens,
    utilizando operações de convolução.}
}

\newglossaryentry{unet}{
  name=U-Net,
  description={Arquitetura de segmentação semântica baseada em
    codificador-decodificador com conexões de salto, amplamente utilizada
    para segmentação de imagens médicas e odontológicas.}
}

\newglossaryentry{gradcam}{
  name=Grad-CAM,
  description={Gradient-weighted Class Activation Mapping: técnica de
    visualização que destaca as regiões da imagem mais relevantes para
    a decisão de uma CNN, gerando mapas de ativação.}
}

\newglossaryentry{gememdigital}{
  name=Gêmeo Digital,
  description={Réplica virtual detalhada de uma estrutura física ou
    biológica, atualizada em tempo real com dados do sensor, usada para
    simulação e planejamento de tratamentos.}
}

\newglossaryentry{sdd}{
  name=SDD (Especificação Dirigida por Software),
  description={Metodologia de desenvolvimento que parte de uma
    especificação formal do sistema para orientar a implementação e
    os testes, garantindo rastreabilidade entre requisitos e código.}
}

\newglossaryentry{tdd}{
  name=TDD (Test-Driven Development),
  description={Metodologia de desenvolvimento onde os testes são escritos
    antes do código de produção, seguindo o ciclo Red-Green-Refactor.}
}

% Glossário como tabela
O glossário completo está disponível no final do livro, na seção
\backmatter. Abaixo listamos os principais termos:

\begin{longtable}{lp{0.65\textwidth}}
  \toprule
  \textbf{Termo} & \textbf{Definição} \\
  \midrule
  \endhead
  Aprendizado Supervisionado & Método de ML treinado com dados rotulados. \\
  Aprendizado Não Supervisionado & Método que descobre padrões sem rótulos. \\
  Deep Learning & Redes neurais com múltiplas camadas. \\
  Grad-CAM & Mapa de ativação para interpretação de CNNs. \\
  Gêmeo Digital & Réplica virtual de sistemas físicos. \\
  Segmentação Semântica & Classificação pixel-a-pixel de imagens. \\
  Transfer Learning & Reutilização de modelos pré-treinados. \\
  U-Net & Arquitetura para segmentação semântica. \\
  \bottomrule
\end{longtable}